\head{定点核查报告：反例、修订与保留事项}
\noindent\textbf{阅读界限。}下列分析由本次译校另行撰写，不是原著的组成部分。所说“原文”均指本次实际查到的上传载体；对于EGA V，尚未核对法文手写原稿，故不能判定错误最早出在作者、英译者还是排版者。未列出的条目不等于已经审定无误。

本报告以编号定位，区分六项有明确反例的断言或推理（E01--E06）、一个另列的证明步骤问题（G01）、成组的文字与符号问题（T01--T10），以及六项保留事项（P01--P06）。这些编号不能简单相加当作“发现了同样多个错误定理”；例如E04、E06批评的是辅助断言或证明判据，而非所服务的全部结论。

\egasection{重点位置总表}
\begin{longtable}{p{.10\textwidth}p{.28\textwidth}p{.48\textwidth}}
\toprule 编号 & 底本位置 & 判定对象\\\midrule\endhead
E01 & IV.1，$0$（14.3.3）；PDF第10页 & 四个条件并不全等价；条件(b)不足。\\
E02 & IV.1，$0$（15.1.4）；PDF第12页 & $K/A$例子漏掉“$f$非单位”。\\
E03 & V合并稿（2.4）；第7页 & 截面相等的判据只假设了一个集合不可约。\\
E04 & V合并稿（2.12）证明；第10页 & “泛线性组合处处非零”的辅助引理错误。\\
E05 & V合并稿（2.15）、（2.18）；第11--12页 & 特征0允许任意态射时，不能要求坏点集本身有限。\\
E06 & V短篇；正文第3页／PDF第4页 & 层同态单射不能替代Jacobian的纤维满秩条件。\\
G01 & V合并稿（2.12）证明；第10页 & 非分歧不等于Zariski局部浸入。\\
T01 & V短篇末两页 & 外幂、秩、商模及行列式对偶出现错置。\\
T02 & V射影丛笔记（1.1）；PS转换第3页，合并稿第87页 & 扭转 $(1)$ 应为 $(n)$，另有遗漏拉回记号。\\
T03 & III.2（6.3.8）；PDF第11页 & 边界模的解消被误说成分量模的解消。\\
T04--T07 & V合并稿§§1--2 & 拉回对象、空集边界、秩不等式、底域等问题。\\
T08--T10 & V短篇定义及微分图 & 域值点、相对下标、维数算式、余法列末项等。\\
\bottomrule
\end{longtable}

\audit{E01}{“等维、等余维且悬链”不充分}{已确认的命题错误；有独立反例及已有文献记录}
\source{\texttt{EGAIV.1.pdf}第10页，原书正文页11，$0_{\mathrm{IV}}$（14.3.3）}
原文把以下条件(b)与“所有极大不可约闭集链等长”(a)等同：$X$等维、等余维、悬链。这里的极大链是不能再向任一端或在中间插入元素的链；不能把“极大”误读成“长度最大的”。

\subsec{一个六点反例}
取集合 $X=\{a,b,c,d,u,v\}$，以如下覆盖关系定义偏序：
\[
 a<c<u,\qquad b<d<v,\qquad a<v,
\]
除传递性迫使的关系外没有其他严格比较关系。规定闭集为向下封闭的子集。下面的线段表示覆盖关系，向上表示偏序增大：
\begin{center}
\begin{tikzpicture}[x=1.7cm,y=1.0cm]
\node (a) at (0,0) {$a$}; \node (b) at (2,0) {$b$};
\node (c) at (0,1) {$c$}; \node (d) at (2,1) {$d$};
\node (u) at (0,2) {$u$}; \node (v) at (2,2) {$v$};
\draw (a)--(c)--(u); \draw (b)--(d)--(v); \draw (a)--(v);
\end{tikzpicture}
\end{center}
这是有限 $T_0$ 空间，因而诺特。每个非空不可约闭集都是一个主向下集
\[
 \downarrow x=\{z:z\leq x\}.
\]
证明如下：每个非空闭集是其有限多个极大元的主向下集之并；若极大元不止一个，它可写成两个真闭子集之并，不可约性便不成立。反之，主向下集包含其唯一极大元，任何闭集分解都必须有一个分量包含该极大元，从而包含整个主向下集。

故不可约闭集链与偏序中的点链一一对应。两个不可约分支是
\[
 C_u=\{a,c,u\},\qquad C_v=\{a,b,d,v\}.
\]
它们的维数均为2，所以 $X$ 等维。最小不可约闭集只有 $\{a\}$、$\{b\}$，且
\[
 \codim(\{a\},X)=2,\qquad \codim(\{b\},X)=2,
\]
前者由 $a<c<u$ 实现，后者由 $b<d<v$ 实现；所以 $X$ 等余维。最后，对每一对可比较端点，其间只有一条饱和链，因此悬链条件成立。

但 $a<c<u$、$b<d<v$、$a<v$ 都是极大链，其长度分别为 $2,2,1$。所以(b)成立而(a)不成立。这个有限拓扑反例已足以反驳原文针对一般诺特 $T_0$ 空间的命题；无需先构造一个环。

\subsec{错误步骤与正确关系}
问题是把\emph{在整个 $X$ 中的余维数}同\emph{在每个不可约分支中的余维数}混为一谈。反例中
\[
 \codim(\{a\},X)=2,
 \qquad \codim(\{a\},C_v)=1,
 \qquad \dim C_v=2.
\]
原条件(b)未保证 $C_v$ 自身等余维。

保留原来的诺特、$T_0$、有限维假设，正确关系是
\[
 (a)\Longleftrightarrow(c)\Longleftrightarrow(d)\Longrightarrow(b),
\]
而(b)的逆向一般为假。可把(b)加强为
\[
 \boxed{X\text{等维且悬链，并且每个不可约分支自身等余维。}}
\]
这一加强条件与(a)等价：在一个不可约分支内，任何极大链都从一个闭点到该分支；悬链性使其长度等于该闭点在该分支内的余维数，该分支的等余维性使这个数等于分支维数，再由全空间等维性得到统一长度。反过来，所有极大链等长显然迫使每个分支的这些链等长，并给出上述三条条件。

原文在（14.3.3）之后用“四条件等价”定义“biéquidimensionnel”；为避免后续歧义，应以(a)的强性质为定义，另将(b)称为弱双等维性。后续（14.3.4）、（14.3.5）不能仅凭弱性质引用。

\textbf{外部交叉核验。}Katharina Heinrich的研究\textit{Some remarks on biequidimensionality of topological spaces and Noetherian schemes}（arXiv:1403.5814，2014）明确指出上述EGA等价关系的问题，并给出包括诺特环谱在内的反例。这里的六点计算是本报告直接列出的验证，不依赖把那篇文章中的图或证明照搬进来。

\audit{E02}{非零元素不一定是零因子}{已确认的遗漏条件}
\source{\texttt{EGAIV.1.pdf}第12页，原书正文页13，（15.1.4）末尾}
原文取不是域的主理想整环 $A$、分式域 $K$，以及 $M=K/A$；正确地指出每个 $f\ne0$ 都满足 $fM=M$，但随后说这些 $f$ 都是 $M$ 上的零因子。

取 $A=\mathbf Z$、$K=\mathbf Q$、$f=1$，乘法映射就是 $M$ 的恒等映射，当然单射；所以该句按字面为假。正确条件为
\[
 0\ne f\in A,\qquad f\notin A^\times.
\]
此时 $1/f\notin A$，所以 $1/f+A$ 是 $K/A$ 的非零元素，却被 $f$ 消去：
\[
 f(1/f+A)=1+A=0.
\]
同时 $fM=M$，因此关联分次模及 $(M/fM)[T]$ 都为零，典范同态为同构。这才给出“准正则不推出正则”的所需反例。无须否定此前（15.1.2）、（15.1.3）的结论。

\audit{E03}{超平面逆像相等不能识别任意闭集}{已确认的遗漏条件}
\source{\texttt{EGAVaiello.pdf}第7页，推论（2.4）的第二句}
原文只假设 $Z$ 不可约，却声称
\[
 Z_\eta=Z'_\eta
 \Longleftrightarrow
 \bigl(f(Z),f(Z')\text{均有限}\bigr)\ \text{或}\ Z=Z'.
\]
取代数闭域 $k$，$X=P=\mathbf P^2_k$、$f=\id$。令 $Z=L$ 为一条直线，取 $p\in P\setminus L$ 为闭点，令 $Z'=L\cup\{p\}$。两个集合均闭，$Z$ 不可约，完全满足所写假设。

泛超平面 $H_\eta$ 不经过固定闭点 $p$：在对偶射影平面中，“经过 $p$”是一个真闭超平面，不包含泛点。于是
\[
 Z'_\eta=(L\cup\{p\})_\eta=L_\eta=Z_\eta.
\]
但 $Z\ne Z'$，并且两个像都包含一条直线，均非有限。这就是反例。

\textbf{可用修正。}若同时要求 $Z'$ 不可约，则可对原文第一句的包含判据应用两次；在非空逆像的情形，两个像都不有限，于是得到 $Z\subseteq Z'$ 和 $Z'\subseteq Z$。本反例只否定第二句的过强结论，不否定第一句。

\audit{E04}{泛线性组合可以有零点}{已确认的辅助断言错误，不是Bertini定理反例}
\source{\texttt{EGAVaiello.pdf}第10页，（2.12）证明末段}
底本使用：非零局部自由模 $E$ 由 $a_1,\ldots,a_n$ 生成，则以代数独立系数形成的 $\sum_iS_i a_i$ 在扩域后处处不消失。

取 $X=\mathbf A^1_k=\Spec k[t]$、$E=\OO_X$，取生成系 $a_1=1,a_2=t$。令 $K=k(S_1,S_2)$。在 $X_K$ 中，
\[
 s=S_1+S_2t
\]
在闭点 $t=-S_1/S_2$ 处消失，因为 $S_2$ 在域 $K$ 中可逆。这个点是扩域后新出现的点，坐标依赖参数。因此“对每一个固定的旧点，其一般线性组合不消失”不能推出“扩域后所有点处都不消失”。

\textbf{修补方向。}不能只修补一个求和下标就保住该论证；应同时考虑截面方程和奇异性条件的关联空间。对光滑浸入 $X\hookrightarrow\mathbf P^n_k$ 的标准Bertini结论，可以用超平面在 $x$ 处的一阶射流条件估算坏超平面的维数，而不是宣称某个泛组合在整个扩域空间上无零点。Stacks Project，Tag 0FD4，Lemma 33.47.3给出了具有明确浸入假设的对应结果及证明。此处仅以它作结论与证明路线的交叉核对；不把其结论不加验证地推广到草稿的所有版本。

\audit{E05}{任意态射下应控制坏点的像}{对特征0无浸入假设扩展的明确反例}
\source{\texttt{EGAVaiello.pdf}第11--12页，（2.14）的旁注、（2.15）、（2.18）}
这里必须区分两个版本。保留非分歧有限型假设时，态射拟有限，有限像与坏点集合有限之间有额外联系；本项不是对那个受限版本的反例。但若按（2.14）旁注及（2.18）的方向，在特征0允许 $f$ 为任意态射，则不能仍把条件写作“$Z$ 有限”。

取 $k=\mathbf C$，
\[
 X=\Spec k[t,\epsilon]/(\epsilon^2),
 \qquad f:X\longrightarrow\mathbf P^2_k
\]
为到固定闭点 $p$ 的常值态射。因为 $X$ 处处带有非零幂零元，它在任何点都不对 $k$ 光滑，所以不光滑点集 $Z$ 就是 $X$，不是有限集。另一方面，泛超平面不经过 $p$，因此
\[
 Y_\eta=f^{-1}(H_\eta)=\varnothing.
\]
空预概形对任何域光滑。故“$Y_\eta$ 光滑当且仅当 $Z$ 有限”在该扩展版本下为假，而“$f(Z)$ 有限”在这个例子中确实成立。

这一例子揭示的是由像到原像时丢掉了拟有限性，绝不反驳（2.12）的点态等价：这里 $Y_\eta$ 没有点，该点态陈述没有产生矛盾。关于（2.17）的“模可分”版本，必须另外核对草稿所引用但未写出的定义，不把不同历史用法混同；但（2.18）只证明点态结果，显然没有说明为何能继续要求坏点集本身有限。

\audit{E06}{余法映射单射不等于光滑}{已确认的证明判据错误}
\source{\texttt{EGAV1-2.ps}正文第3页／转换PDF第4页}
原稿称某个微分正合列的第一个映射单射，恰好等价于零点概形光滑且维数正确。即使把误写的相对底 $X$ 改成 $S$，只说\emph{层映射单射}也不够。

取特征0的域 $k$，令
\[
 X=\Spec k[x,y],\qquad
 W=V(y^2-x^3),\qquad B=k[x,y]/(y^2-x^3).
\]
因 $y^2-x^3$ 是非零因子，余法模 $J/J^2\simeq B$，余法映射为
\[
 B\longrightarrow B\,dx\oplus B\,dy,
 \qquad h\longmapsto(-3x^2h,2yh).
\]
$B$ 是整环，且 $2y\ne0$，所以这个映射单射。但在原点 $\mm=(x,y)$ 的剩余域上，它变为零映射
\[
 k\longrightarrow k^2.
\]
$W$ 在原点的局部环维数为1，而 $\dim_k\mm/\mm^2=2$，不是正则局部环，因此不是光滑点。

正确Jacobian条件应以\emph{纤维上的秩}表述：在光滑环境中，用 $r$ 个方程切出预期余维 $r$ 的光滑零点概形，需要这些方程的微分在该点的余切空间中线性无关。这里不独立断言草稿命题7的结论错误；所否定的是支持它的一句过强判据。

\audit{G01}{非分歧态射不必Zariski局部是浸入}{原证明缺少合法的局部化步骤}
\source{\texttt{EGAVaiello.pdf}第10页，（2.12）的逆向证明开头}
令 $\operatorname{char}k\ne2$，取
\[
 f:\mathbf G_{m,k}\longrightarrow\mathbf G_{m,k},\qquad u=t^2.
\]
对应环扩张 $k[u,u^{-1}]\to k[t,t^{-1}]$ 有基 $1,t$，为秩2的有限自由扩张；多项式 $T^2-u$ 的导数 $2T$ 在这里可逆，所以态射 étale，特别非分歧。

然而不能通过缩小源、靶的非空Zariski开集就把它变成浸入。这样的限制仍在函数域上给出次数2的扩张 $k(u)\subset k(t)$；若它是占优浸入，其像在不可约靶中稠密而局部闭，因而在某个开集上为同构，函数域次数必须为1，矛盾。

因此，“$f$非分歧，于是把 $X$ 直接写成 $P$ 的某个Zariski开集内由方程切出的子概形”这一步不合法。可以考虑 étale 局部浸入或直接的切空间论证，但必须写出相应变换及性质的回降；本报告不把这一省略用“显然”填掉。

\audit{T01}{短篇末尾的外幂与行列式类型不一致}{已定位的转录／排印问题}
\source{\texttt{EGAV1-2.ps}转换PDF第5--6页，正文第4--5页}
本项集中列出同一段计算的多处错置，不把每个错字算作一个错误定理。设 $Q\to R$ 为秩分别为 $q,r$ 的局部自由模同态。表示非满射点的标准截面属于
\[
 \mathscr Hom\left(\bigwedge^r Q,\bigwedge^r R\right),
\]
其零点由全部 $r$ 阶子式定义；不能把底本渲染中的 $\Lambda^1Q\to\Lambda\sqrt M$ 当作有意义的公式。

在其后 $M$ 秩为 $m$、子模 $P,Q$ 秩分别为 $p,q$ 的段落中，余秩1情形所需的是
\[
 p+q=m+1,\qquad S=M/Q,\qquad R=M/P.
\]
底本却写 $p+1=m+1$ 与 $M/W$，其中 $W$ 是底空间。再令
\[
 D=\det M\otimes(\det P)^{-1}\otimes(\det Q)^{-1}.
\]
截面 $a\in\Gamma(P\otimes D)$ 对应的同态是 $D^{-1}\to P$，所以
\[
 L=D^{-1}=\det P\otimes\det Q\otimes(\det M)^{-1},
\]
而不是底本写的末因子 $\det M$。这只校正类型及对偶方向；原文件缺失的最后交换图仍不能由这些计算冒充恢复。

\audit{T02}{射影丛可逆层定理的扭转下标}{已确认的公式误写}
\source{\texttt{EGAV6.ps}转换PDF第3页；\texttt{EGAVaiello.pdf}第87页，（1.1）}
\textbf{被检查段落的完整陈述摘译。}设 $S$ 为预概形，$E$ 为 $S$ 上有限秩局部自由模，且每点处秩至少为2，$P=\mathbf P(E)$ 为它定义的射影丛。对 $P$ 上每个可逆模 $L$，可以找到以 $\mathbf Z$ 为指标的一族两两不交、覆盖 $S$ 的开集 $(S_n)$，以及 $S$ 上的可逆模 $M$，使 $L$ 在 $f^{-1}(S_n)$ 上的限制同构于原稿写作 $M\otimes\OO_P(1)=f^*(M)(n)$ 的限制。族 $(S_n)$ 被唯一确定，而 $M$“在唯一同构意义下确定”。

正确的扭转式应为
\[
 f^*M\otimes\OO_P(n)=f^*(M)(n).
\]
这里既要补出从 $S$ 到 $P$ 的拉回，也要把第一个 $(1)$ 改为 $(n)$。取 $S=\Spec k$、$P=\mathbf P^1_k$、$L=\OO_P(2)$，则 $M$ 拉回为平凡线丛，$\OO(2)$ 不同构于 $\OO(1)$；例如两者的全局截面维数分别为3与2。这是公式笔误，不是否定正确的Picard群分解。关于“唯一同构”的严格含义，见P06。

\audit{T03}{边界模的解消与分量模的解消混写}{原文中可定位的对象误写}
\source{\texttt{EGAIII.2.pdf}第11页，原书正文页12，（6.3.8）}
底本先定义 $X^B_{i,\bullet}$ 是 $B_i(P_\bullet)$ 的投射解消，随后却说 $X^B_{i,\bullet}\otimes_AA'$ 一般不再是 $P_i\otimes_AA'$ 的解消。它在变基前就不是被指定为 $P_i$ 的解消。应核对为边界对象 $B_i(P_\bullet)\otimes_AA'$，并明确它到 $B_i(P_\bullet\otimes_AA')$ 的典范映射。非平坦变基不保证后两个边界对象相同。此项不据此否定最终构造的Tor变基同态。

\audit{T04}{拉回层的承载空间与箭头靶错置}{对象类型错误}
\source{\texttt{EGAVaiello.pdf}第6页及第9页}
第1节把 $\G$ 写成 $X$ 上层 $\F$ 在 $H$ 上的拉回，但已定义的关联关系是 $Y=X\times_PH$，存在 $Y\to X$，并没有一般的典范 $H\to X$。应为拉回到 $Y$。（2.8）又写由 $Y_\eta\to Y$ 的正则性得 $\OO_{X,x}\to\OO_{Y_\eta,y}$ 正则；证明所需的态射是复合 $Y_\eta\to X$，须指出其性质，不能混同这两个箭头。

\audit{T05}{“有限”与“维数等于0”遗漏空集}{边界约定未统一}
\source{\texttt{EGAVaiello.pdf}第9页，（2.9）、（2.11）}
坏点集可以为空。若按EGA的约定 $\dim\varnothing=-\infty$，则“$f(Z)$有限”不等价于“$\dim\overline{f(Z)}=0$”，因为 $Z=\varnothing$ 即反例。在相应可构造性条件下，应改为维数 $\leq0$，或明确分出空集。这是边界表述问题，不是一个非空奇异几何反例。

\audit{T06}{方程个数与商模秩的下标}{公式及不等式误写}
\source{\texttt{EGAVaiello.pdf}第10页，（2.12）证明}
已列方程为 $f_1,\ldots,f_p$，故商模应除以 $\sum_{i=1}^{p}\OO_Udf_i$，不是求和到 $n$。在这些微分独立时，若 $\Omega^1_{U/k}$ 秩为 $n$，商的秩为 $n-p$，所以商非零要求 $p<n$，不是 $p\leq n$。修正这两点仍不能挽救E04的错误辅助引理。

\audit{T07}{局部环的底域不是其剩余域}{代数结构未定义／底域误写}
\source{\texttt{EGAVaiello.pdf}第11页，（2.14）}
$Y_\eta$ 是 $k(\eta)$-预概形，因此 $\OO_{Y_\eta,y}$ 自然是 $k(\eta)$-代数。它有到 $k(y)$ 的商同态，但一般没有反向的典范 $k(y)\to\OO_{Y_\eta,y}$。故所写
\[
 P(\OO_{Y_\eta,y},k(y))
\]
没有给出所需的底域代数结构；相邻陈述及证明均指向 $k(\eta)$。改底域是必要的类型修订，不意味着该段所有其余推理都已证明。

\audit{T08}{短篇中的测试域、次数与相对下标}{同一载体内的符号不一致}
\source{\texttt{EGAV1-2.ps}转换PDF第2--5页}
定义1的二次初始项应落在 $\Sym^2(\mm_x/\mm_x^2)$。定义3既量化扩域 $k'$，相应值点也应以 $k'$ 为值域，不能写 $k$。命题5的相对底先为 $Y$，后误写尚未引入的 $S$；分析 $\mu$ 时又把 $X/Y$ 换成 $X/S$。最后“$M$满射”“$V$满射”对应图中的 $\mu,\nu$。这些错置在译文均明示，不把数学对象悄悄换掉。

\audit{T09}{相对维数算式错误}{可直接验算的排印错误}
\source{\texttt{EGAV1-2.ps}转换PDF第4页}
原式 $(m+1)-(d+1)=m-1$ 的左边等于 $m-d$。因 $X/Y$ 相对维数为 $d$、$Y/S$ 为 $m$，在 $X$ 内切出 $d+1$ 个独立条件所对应的计算应是 $(m+d)-(d+1)=m-1$。

\audit{T10}{余法正合列的末项与解释层次}{类型错误；另有E06所述实质问题}
\source{\texttt{EGAV1-2.ps}转换PDF第4页}
若 $W\hookrightarrow X$ 的理想为 $J$，则相对于 $S$ 的余法列是
\[
 J/J^2\longrightarrow\Omega^1_{X/S}\otimes\OO_W
      \longrightarrow\Omega^1_{W/S}\longrightarrow0.
\]
底本把末项误写为 $\Omega^1_{X/S}$。若 $W$ 由局部自由模 $F$ 的截面定义，且把整个列作为 $W$ 上的模列，则首项写为 $F^\vee|_W$。但若统一视作 $X$ 上的模层，首项 $F^\vee$ 也可先经限制映射再进入该列；\textbf{这时首项本身不另构成错误}。真正明确误写的是余法列末项。无论采用哪一种解释，仍必须区分层单射与纤维满秩，不能跳过E06。

\audit{P01}{原稿本来没有写完的内容}{缺文，不补造}
\source{V合并稿（2.10）、（2.16）；V短篇“命题6”及末尾}
（2.10）只有一句拟写推论的提示，没有正式条件和证明；（2.16）的英译编者直接注明某句无法理解；短篇说应归纳成命题6，却未写出编号陈述；文件末尾则在承诺一个交换图后结束。这些是载体实际的内容边界。既不能说相应证明已译全，也不能因为原作者预期“应当容易”就补写为原文。

\audit{P02}{只知上同调层凝聚，能否直接使用逐项拟凝聚的计算}{证明适用范围待补核；不判定命题为假}
\source{\texttt{EGAIII.2.pdf}第7页，（6.2.5）；对照第6页（6.2.2）、（6.2.3）}
（6.2.5）的陈述只假设 $\cH^i(\K^\bullet)$ 凝聚，没有明写每个 $\K^i$ 拟凝聚；证明却直接援引对逐项拟凝聚复形建立的仿射Čech计算，并使用该类上的截面函子正合性。仅从“上同调层凝聚”不能直接推得“复形每项拟凝聚”。

这里需要另给合法替换或更一般的比较定理。作为\emph{另加条件后的论证}，若 $\K^\bullet$ 有下界，则可从第二超上同调谱序列
\[
 E_2^{p,q}=R^pf_*\cH^q(\K^\bullet)\Longrightarrow
 \mathbb H^{p+q}(f,\K^\bullet)
\]
以及本征直像的凝聚性逐个总次数处理：$p\geq0$、$q$有统一下界，使固定总次数上只剩有限项。若逐项拟凝聚，则原稿的有限仿射覆盖路线也有明确适用条件。但这些补充不能冒充原文已经写出的普遍证明；一般无界情形及原著第0章的完整约定仍需进一步对照。

\audit{P03}{一般底上相对除子函子的参数空间}{记号与约定待核}
\source{\texttt{EGAVaiello.pdf}第5页}
底本在一般预概形 $S$ 上写 $\mathbf P(\Sym_k^m(E^\vee))$，而 $k$ 并未作为一般底引入；还涉及射影化采用商线还是子线约定、相对维数0情形、空除子的次数0分支。不能仅凭熟悉的射影空间超曲面参数化就把这句话无条件认定为已经完全严谨。至少需要统一为相对于 $\OO_S$ 的构造并核对全部约定。

\audit{P04}{“正规”与“几何正规”的不对称表述}{待核；不把文字猜测升格为原作者本意}
\source{\texttt{EGAVaiello.pdf}第11页，（2.15）}
坏点集的最后一项只写“非正规”，结论却写“几何正规”。若不限定完美底域，两者不能任意互换。上下文可能漏掉一个“几何”，也可能另有限制；没有核对手稿前，应保留原句并明确疑点。关于有限性假设的另一个问题，已经在E05给出反例，应与本项分开。

\audit{P05}{特征2的二次型退化约定}{术语未明确，不能贸然认定数学命题错误}
\source{\texttt{EGAV1-2.ps}定义1及Hessian行列式段}
若二次型“非退化”被定义为极化双线性型非退化，那么后文的行列式判据与这个定义相容。但若想用“对应的射影二次超曲面光滑”作为非退化的含义，特征2时必须小心。例如取特征2的域，
\[
 q(x,y,z)=x^2+yz.
\]
其极化双线性型在三维空间上有非零根空间（$x$方向），故行列式为零；然而射影曲线 $q=0$ 光滑，因为偏导是 $(0,z,y)$，偏导全零只能有 $y=z=0$，而此时 $q=x^2=0$，没有射影点。因此仅写“退化”不足以决定想采用哪一类判据。本报告保留为约定待核，而不是宣布原稿在所有特征下的定义错误。

\audit{P06}{可逆模“在唯一同构意义下确定”}{必须说明保留了哪些结构}
\source{\texttt{EGAV6.ps}（1.1）、（1.3）；合并稿第87页}
若只给可逆模的同构类，通常不能说两个实现之间只有一个同构：例如 $\Spec k$ 上一维向量空间 $M$ 的自同构群是 $k^\times$。若同时保留一个指定的同构 $f^*M(n)\simeq L$，并要求比较同构与它相容，则“唯一”可能有正确的范畴意义。应区别“$M$的同构类唯一”与“带指定比较同构的数据在唯一同构意义下唯一”。没有固定数据的裸说法不够严格；本项不反驳正确的Picard群分类。

\egasection{不计入数学反例的边界与识读事项}
EGA IV（14.3.4.1）的原扫描式是多个等号；文本抽取曾把第一个等号读成减号。本译稿按页面图像恢复等号，\textbf{不把识别错误记成原著错误}。同样，EGA V 的像闭包上的横线若在文本抽取中丢失，也要按页面恢复。

（14.3.5.1）写“每个闭子集 $Y$”时，如果容许 $Y=\varnothing$，其左侧按原文约定会出现 $-\infty+(+\infty)$；宜明确限定非空 $Y$，或另说明空集的约定。这是边界记号提醒，未与E01的实质反例混算。

\egasection{本次没有得出的结论}
没有完成本次九文件的全译，没有逐行验证全部EGA III、IV证明，也没有证明所有未标注命题正确。没有以EGA V英译的乱码断定Grothendieck手稿必然有同样文字错误。没有因E04或G01而宣布Bertini定理错误。没有把PS格式转换、正文抽取或编译成功计为翻译完成。
